\problemname{Duplantis}
\noindent

Duplantis har ett problem. Han kan inte bestämma sig för om han ska slå så många världsrekord som möjligt eller ett 
världsrekord som är så bra som möjligt.

Just nu är hans världsrekord på $v$ cm. Duplantis har förutspått hur hans form kommer vara de resterande $N$ tävlingarna 
i karriären, och har räknat fram de positiva heltalen $a_1, a_2, \dots, a_N$. Talet $a_i$ innebär att Duplantis kan 
hoppa mellan $0$ och $a_i$ cm i den $i$:te tävlingen. Han kan bestämma fritt hur högt han kommer hoppa i en viss tävling,
så länge som höjden är ett heltal mellan $0$ och $a_i$ cm.

För att slå ett världsrekord måste Duplantis hoppa strikt högre än sitt förra världsrekord. Det är bara möjligt att slå som mest 
ett världsrekord i en och samma tävling. I stavhopp så hoppar man alltid ett helt antal centimeter.

Din uppgift är att räkna fram $r$, det maximala antalet fler världsrekord Duplantis kan slå. 
Du ska också räkna ut talen $b_1, b_2, \dots, b_r$, där $b_i$ är det maximala möjliga slutgiltiga världsrekordet 
givet att Duplantis slår exakt $i$ fler världsrekord.

\section*{Indata}
Den första raden innehåller två heltal $N$ och $v$ ($1 \leq N \leq 3 \cdot 10^5$, $0 \leq v \leq 10^9$), där $N$ är antalet tävlingar och $v$ 
är Duplantis nuvarande världsrekord.

Den andra raden innehåller $N$ heltal $a_1, a_2, \dots, a_N$ ($0 \leq a_i \leq 10^9$). 

\section*{Utdata}
På den första raden, skriv ut $r$, det maximala antalet fler världsrekord Duplantis kan slå.
På den andra raden, skriv ut de $r$ talen $b_1, b_2, \dots, b_r$, där $b_i$ är det maximala möjliga slutgiltiga världsrekordet om Duplantis 
slår exakt $i$ fler världsrekord.

Notera att $r$ kan vara noll, i så fall ska du lämna den andra raden tom (att inte skriva ut en tom rad är också OK).

\section*{Poängsättning}
Din lösning kommer att testas på en mängd testfallsgrupper.
För att få poäng för en grupp så måste du klara alla testfall i gruppen.

\noindent
\begin{tabular}{| l | l | p{12cm} |}
  \hline
  \textbf{Grupp} & \textbf{Poäng} & \textbf{Gränser} \\ \hline
  $1$    & $20$       & $v < a_1 < a_2 < \dots < a_N$ \\ \hline
  $2$    & $20$        & $a_1 > a_2 > \dots > a_N$ \\ \hline
  $3$    & $17$        & $N = 2$ \\ \hline
  $4$    & $43$       & Inga ytterligare begränsningar. \\ \hline
\end{tabular}
